% =====================================================================
% Paper 07-ext — SO(10) anomaly cancellation auxiliary note
%   (HEALTHIEST of Wave 4/5; aligned with Math319 V1 verification target)
% Status: [DRAFT] [NEEDS_UPDATE] — Wave 4 anomaly sector
% AUDIT-FLAG: AUDIT-2026-05-02-Wave7-Aux-Epoch-Overclaim
%   (Math314-AddC extension covering Wave 4/5 papers Paper-06..08).
% Rev 4 (2026-05-06):
%   (a) RHN physics correction. The earlier draft asserted that the
%       right-handed neutrino (RHN) singlet 1(5) of SO(10), with
%       SU(5)xU(1)_chi quantum numbers, "cancels the U(1)_Y anomaly of
%       the 5̄(-3) representation". This is incorrect: the RHN has
%       hypercharge Y = 0 in the SM (Y = T_3R + (B-L)/2 = -1/2 + 1/2
%       = 0), and therefore contributes ZERO to any U(1)_Y trace
%       Tr(Y^n). The 5̄+10 of SU(5) already saturates all SM gauge
%       anomaly cancellations (Tr(U(1)_Y^3), Tr(SU(2)_L^3),
%       Tr(SU(3)_c^3), Tr(U(1)_Y · SU(2)_L^2),
%       Tr(U(1)_Y · SU(3)_c^2), Tr(U(1)_Y)) per generation. The RHN's
%       genuine roles are: (i) closure of the Witten SU(2) global
%       anomaly when the chiral SU(2) doublet count must be even;
%       (ii) saturation of the mixed gravitational—U(1)_{B-L} anomaly
%       (Tr(B-L)) when (B-L) is gauged in the SO(10) extension;
%       (iii) completion of the 16 spinor of SO(10) so that the SO(10)
%       gauge group itself is anomaly-free at the algebraic level
%       (SO(10) is anomaly-free as a simple Lie algebra; 16 is the
%       spinor that realises this).
%   (b) Anomaly list precision. Six SM anomaly coefficients are now
%       split into perturbative gauge anomalies (chiral triangle
%       diagrams) and the gravitational mixed anomaly Tr(U(1)_Y) (the
%       graviton-graviton-U(1)_Y triangle). The Witten SU(2) global
%       anomaly is a separate, non-perturbative, mod-2 statement and
%       is not an entry in the perturbative six-coefficient list.
%   (c) "Numerically to six-digit precision" replaced by "symbolically
%       exact zero" throughout (the cancellation is symbolic, not a
%       numerical estimate).
%   (d) Explicit decoupling from Paper-07's three-generation theorem
%       (which does not stand per Math314-AddC); this paper is an
%       independent anomaly-cancellation auxiliary note. Three
%       generations enter ONLY via the Tr(...) being a per-generation
%       statement that holds for any number of generations N_g.
% Compile: bash _shared/build-paper.sh Paper-07-ext-SO10-Anomaly/Paper-07-ext
% Canonical archive: Math157, Math157-AddD, Math319 (V1 target),
%                    Math314-AddC
% =====================================================================
% --- TECT canonical preamble (see _shared/tect-preamble.tex) ----------
% =====================================================================
% TECT canonical LaTeX preamble (revtex4-2 / single-column / theorem-safe)
% =====================================================================
% Maintainer: Jusang Lee (jtkor@outlook.com)
% Binding from: 2026-05-06
% Purpose: a single source-of-truth preamble for every TECT paper
% (Paper-00 .. Paper-10 + Auxiliary notes). Designed to (a) eliminate
% font-corruption on pdflatex (T1 fontenc + lmodern), (b) make
% theorem-heavy mathematical-physics content render cleanly in a
% single-column layout (PRD class option, NOT PRL two-column), (c)
% provide consistent theorem numbering across all TECT papers via
% amsthm with a shared counter set, (d) make hyperref + cleveref
% cross-references resolve cleanly with the standard two-pass
% pdflatex compile.
%
% Usage in each Paper-NN.tex:
%   % =====================================================================
% TECT canonical LaTeX preamble (revtex4-2 / single-column / theorem-safe)
% =====================================================================
% Maintainer: Jusang Lee (jtkor@outlook.com)
% Binding from: 2026-05-06
% Purpose: a single source-of-truth preamble for every TECT paper
% (Paper-00 .. Paper-10 + Auxiliary notes). Designed to (a) eliminate
% font-corruption on pdflatex (T1 fontenc + lmodern), (b) make
% theorem-heavy mathematical-physics content render cleanly in a
% single-column layout (PRD class option, NOT PRL two-column), (c)
% provide consistent theorem numbering across all TECT papers via
% amsthm with a shared counter set, (d) make hyperref + cleveref
% cross-references resolve cleanly with the standard two-pass
% pdflatex compile.
%
% Usage in each Paper-NN.tex:
%   % =====================================================================
% TECT canonical LaTeX preamble (revtex4-2 / single-column / theorem-safe)
% =====================================================================
% Maintainer: Jusang Lee (jtkor@outlook.com)
% Binding from: 2026-05-06
% Purpose: a single source-of-truth preamble for every TECT paper
% (Paper-00 .. Paper-10 + Auxiliary notes). Designed to (a) eliminate
% font-corruption on pdflatex (T1 fontenc + lmodern), (b) make
% theorem-heavy mathematical-physics content render cleanly in a
% single-column layout (PRD class option, NOT PRL two-column), (c)
% provide consistent theorem numbering across all TECT papers via
% amsthm with a shared counter set, (d) make hyperref + cleveref
% cross-references resolve cleanly with the standard two-pass
% pdflatex compile.
%
% Usage in each Paper-NN.tex:
%   \input{../_shared/tect-preamble.tex}
%   \begin{document}
%   ...
%   \end{document}
%
% Build (every paper): `pdflatex Paper-NN && pdflatex Paper-NN`
% (the second pass resolves \ref{} / \cref{} forward references).
% A bash/PowerShell wrapper that automates this is at
% Docs/papers/papers/_shared/build-paper.sh / build-paper.ps1.
% =====================================================================

% ---- Document class --------------------------------------------------
% PRD class option of revtex4-2 with [reprint,onecolumn] gives the same
% APS heading / by-line / abstract style as PRL but in a wider single
% column that handles long display equations and theorem environments
% without column-break artefacts. nofootinbib keeps citations out of
% footnotes. The 11pt option restores standard font metrics; with T1
% fontenc + lmodern this gives non-bitmap glyphs.
\documentclass[%
  aps,prd,reprint,onecolumn,nofootinbib,11pt,%
  amsmath,amssymb,floatfix,longbibliography%
]{revtex4-2}

% ---- Encoding & fonts ------------------------------------------------
\usepackage[T1]{fontenc}        % proper 8-bit font encoding (no font corruption)
\usepackage[utf8]{inputenc}     % allow UTF-8 source
\usepackage{lmodern}            % scalable Latin Modern (replaces bitmap CM)
\usepackage{textcomp}           % extra text symbols
\usepackage{microtype}          % subtle kerning improvements

% ---- UTF-8 → LaTeX character mappings --------------------------------
% Maps the Unicode characters that occur in TECT body text (legacy from
% the chat-content auto-archival rule, CLAUDE.md §4) to LaTeX-safe
% commands. Without these declarations, T1+inputenc rejects the chars
% with "Unicode character X not set up for use with LaTeX".
% Adding declarations centrally (here) is preferable to per-file edits.
\DeclareUnicodeCharacter{00A7}{\S}              % §  (section sign)
\DeclareUnicodeCharacter{00B2}{\ensuremath{^{2}}}        % ²
\DeclareUnicodeCharacter{00B3}{\ensuremath{^{3}}}        % ³
\DeclareUnicodeCharacter{00B5}{\ensuremath{\mu}}         % µ (micro)
\DeclareUnicodeCharacter{00B7}{\ensuremath{\cdot}}       % · (middle dot)
\DeclareUnicodeCharacter{00D7}{\ensuremath{\times}}      % ×
\DeclareUnicodeCharacter{00E4}{\"a}             % ä
\DeclareUnicodeCharacter{00E9}{\'e}             % é
\DeclareUnicodeCharacter{00F6}{\"o}             % ö
\DeclareUnicodeCharacter{00FC}{\"u}             % ü
\DeclareUnicodeCharacter{010C}{\v{C}}           % Č
\DeclareUnicodeCharacter{010D}{\v{c}}           % č
\DeclareUnicodeCharacter{0394}{\ensuremath{\Delta}}      % Δ
\DeclareUnicodeCharacter{03A3}{\ensuremath{\Sigma}}      % Σ
\DeclareUnicodeCharacter{03A9}{\ensuremath{\Omega}}      % Ω
\DeclareUnicodeCharacter{03B1}{\ensuremath{\alpha}}      % α
\DeclareUnicodeCharacter{03B2}{\ensuremath{\beta}}       % β
\DeclareUnicodeCharacter{03B3}{\ensuremath{\gamma}}      % γ
\DeclareUnicodeCharacter{03B4}{\ensuremath{\delta}}      % δ
\DeclareUnicodeCharacter{03B5}{\ensuremath{\varepsilon}} % ε
\DeclareUnicodeCharacter{03BA}{\ensuremath{\kappa}}      % κ
\DeclareUnicodeCharacter{03BB}{\ensuremath{\lambda}}     % λ
\DeclareUnicodeCharacter{03BC}{\ensuremath{\mu}}         % μ
\DeclareUnicodeCharacter{03BD}{\ensuremath{\nu}}         % ν
\DeclareUnicodeCharacter{03C0}{\ensuremath{\pi}}         % π
\DeclareUnicodeCharacter{03C1}{\ensuremath{\rho}}        % ρ
\DeclareUnicodeCharacter{03C3}{\ensuremath{\sigma}}      % σ
\DeclareUnicodeCharacter{03C4}{\ensuremath{\tau}}        % τ
\DeclareUnicodeCharacter{03C6}{\ensuremath{\varphi}}     % φ
\DeclareUnicodeCharacter{03C7}{\ensuremath{\chi}}        % χ
\DeclareUnicodeCharacter{03C8}{\ensuremath{\psi}}        % ψ
\DeclareUnicodeCharacter{03C9}{\ensuremath{\omega}}      % ω
\DeclareUnicodeCharacter{2013}{--}              % – (en dash)
\DeclareUnicodeCharacter{2014}{---}             % — (em dash)
\DeclareUnicodeCharacter{2018}{`}               % ‘
\DeclareUnicodeCharacter{2019}{'}               % ’
\DeclareUnicodeCharacter{201C}{``}              % “
\DeclareUnicodeCharacter{201D}{''}              % ”
\DeclareUnicodeCharacter{2070}{\ensuremath{^{0}}}        % ⁰
\DeclareUnicodeCharacter{2071}{\ensuremath{^{i}}}        % ⁱ
\DeclareUnicodeCharacter{2122}{\textsuperscript{TM}}     % ™
\DeclareUnicodeCharacter{2126}{\ensuremath{\Omega}}      % Ω (ohm sign)
\DeclareUnicodeCharacter{210F}{\ensuremath{\hbar}}       % ℏ
\DeclareUnicodeCharacter{2115}{\ensuremath{\mathbb{N}}}  % ℕ
\DeclareUnicodeCharacter{2102}{\ensuremath{\mathbb{C}}}  % ℂ
\DeclareUnicodeCharacter{211D}{\ensuremath{\mathbb{R}}}  % ℝ
\DeclareUnicodeCharacter{2124}{\ensuremath{\mathbb{Z}}}  % ℤ
\DeclareUnicodeCharacter{211A}{\ensuremath{\mathbb{Q}}}  % ℚ
\DeclareUnicodeCharacter{2192}{\ensuremath{\to}}         % →
\DeclareUnicodeCharacter{21D2}{\ensuremath{\Rightarrow}} % ⇒
\DeclareUnicodeCharacter{2200}{\ensuremath{\forall}}     % ∀
\DeclareUnicodeCharacter{2203}{\ensuremath{\exists}}     % ∃
\DeclareUnicodeCharacter{2208}{\ensuremath{\in}}         % ∈
\DeclareUnicodeCharacter{2209}{\ensuremath{\notin}}      % ∉
\DeclareUnicodeCharacter{2227}{\ensuremath{\wedge}}      % ∧
\DeclareUnicodeCharacter{2228}{\ensuremath{\vee}}        % ∨
\DeclareUnicodeCharacter{2229}{\ensuremath{\cap}}        % ∩
\DeclareUnicodeCharacter{222A}{\ensuremath{\cup}}        % ∪
\DeclareUnicodeCharacter{2245}{\ensuremath{\cong}}       % ≅
\DeclareUnicodeCharacter{2248}{\ensuremath{\approx}}     % ≈
\DeclareUnicodeCharacter{2260}{\ensuremath{\neq}}        % ≠
\DeclareUnicodeCharacter{2264}{\ensuremath{\leq}}        % ≤
\DeclareUnicodeCharacter{2265}{\ensuremath{\geq}}        % ≥
\DeclareUnicodeCharacter{22C5}{\ensuremath{\cdot}}       % ⋅
\DeclareUnicodeCharacter{2295}{\ensuremath{\oplus}}      % ⊕
\DeclareUnicodeCharacter{2297}{\ensuremath{\otimes}}     % ⊗

% ---- Mathematics & symbols ------------------------------------------
\usepackage{amsmath,amssymb,bm,mathrsfs,mathtools}
\usepackage{amsthm}             % theorem environments (load BEFORE hyperref)

% ---- Theorem environments (shared counter set, TECT canonical) -------
% All theorems / lemmas / propositions / corollaries / remarks share a
% single counter so cross-references read consistently across a paper.
% Definitions get a separate counter (definitions are numbered in their
% own sequence by editorial convention).
\newtheorem{theorem}{Theorem}
\newtheorem{lemma}[theorem]{Lemma}
\newtheorem{proposition}[theorem]{Proposition}
\newtheorem{corollary}[theorem]{Corollary}

\theoremstyle{remark}
\newtheorem{remark}[theorem]{Remark}
\newtheorem{example}[theorem]{Example}

\theoremstyle{definition}
\newtheorem{definition}[theorem]{Definition}
\newtheorem{conjecture}[theorem]{Conjecture}
\newtheorem{hypothesis}[theorem]{Hypothesis}

% ---- Tables / floats / graphics --------------------------------------
\usepackage{booktabs}           % \toprule / \midrule / \bottomrule
\usepackage{array}
\usepackage{graphicx}
\usepackage{enumitem}           % \begin{itemize}[leftmargin=...] options
\usepackage{hyphenat}           % \hyp{} for breakable hyphens in \texttt
\usepackage{seqsplit}           % \seqsplit{} for long unbreakable strings
\sloppy                         % allow stretchier inter-word spacing to
                                % reduce overfull \hbox warnings on long URLs
                                % and long \texttt tokens (paragraph-level).
\emergencystretch=3em           % last-resort hbox stretch budget

% ---- Cross-references (hyperref last among the standard packages,
%      cleveref AFTER hyperref) ----------------------------------------
\usepackage[%
  colorlinks=true,%
  linkcolor=blue,%
  citecolor=blue,%
  urlcolor=blue,%
  breaklinks=true,%
  bookmarks=true,%
  bookmarksnumbered=true%
]{hyperref}
\usepackage[capitalise,nameinlink]{cleveref}

% ---- TECT-canonical author block macros -----------------------------
% (defined here so every paper has the same author / affiliation style)
\newcommand{\tectauthor}{%
  \author{Jusang Lee}%
  \email{jtkor@outlook.com}%
  \affiliation{Independent researcher, Republic of Korea}%
}

% ---- TECT-canonical archive footer ----------------------------------
\newcommand{\tectarchivefooter}{%
  \par\medskip
  \noindent\small\textit{Canonical archive}:
  \url{https://github.com/TECT-OpenPhysics/TECT}.\par%
}

% ---- TECT TODO / AUDIT-FLAG / HONEST-SCOPE inline marks --------------
% Used in drafts to highlight pending audit items without disturbing
% the body text style. Suppress via \tectsuppressmarks (define empty).
\providecommand{\tectauditflag}[1]{\textcolor{red}{\textbf{[AUDIT-FLAG: #1]}}}
\providecommand{\tecttodo}[1]{\textcolor{orange}{\textbf{[TODO: #1]}}}
\providecommand{\tecthonestscope}[1]{\textit{Honest scope: #1.}}

% =====================================================================
% End of TECT canonical preamble
% =====================================================================

%   \begin{document}
%   ...
%   \end{document}
%
% Build (every paper): `pdflatex Paper-NN && pdflatex Paper-NN`
% (the second pass resolves \ref{} / \cref{} forward references).
% A bash/PowerShell wrapper that automates this is at
% Docs/papers/papers/_shared/build-paper.sh / build-paper.ps1.
% =====================================================================

% ---- Document class --------------------------------------------------
% PRD class option of revtex4-2 with [reprint,onecolumn] gives the same
% APS heading / by-line / abstract style as PRL but in a wider single
% column that handles long display equations and theorem environments
% without column-break artefacts. nofootinbib keeps citations out of
% footnotes. The 11pt option restores standard font metrics; with T1
% fontenc + lmodern this gives non-bitmap glyphs.
\documentclass[%
  aps,prd,reprint,onecolumn,nofootinbib,11pt,%
  amsmath,amssymb,floatfix,longbibliography%
]{revtex4-2}

% ---- Encoding & fonts ------------------------------------------------
\usepackage[T1]{fontenc}        % proper 8-bit font encoding (no font corruption)
\usepackage[utf8]{inputenc}     % allow UTF-8 source
\usepackage{lmodern}            % scalable Latin Modern (replaces bitmap CM)
\usepackage{textcomp}           % extra text symbols
\usepackage{microtype}          % subtle kerning improvements

% ---- UTF-8 → LaTeX character mappings --------------------------------
% Maps the Unicode characters that occur in TECT body text (legacy from
% the chat-content auto-archival rule, CLAUDE.md §4) to LaTeX-safe
% commands. Without these declarations, T1+inputenc rejects the chars
% with "Unicode character X not set up for use with LaTeX".
% Adding declarations centrally (here) is preferable to per-file edits.
\DeclareUnicodeCharacter{00A7}{\S}              % §  (section sign)
\DeclareUnicodeCharacter{00B2}{\ensuremath{^{2}}}        % ²
\DeclareUnicodeCharacter{00B3}{\ensuremath{^{3}}}        % ³
\DeclareUnicodeCharacter{00B5}{\ensuremath{\mu}}         % µ (micro)
\DeclareUnicodeCharacter{00B7}{\ensuremath{\cdot}}       % · (middle dot)
\DeclareUnicodeCharacter{00D7}{\ensuremath{\times}}      % ×
\DeclareUnicodeCharacter{00E4}{\"a}             % ä
\DeclareUnicodeCharacter{00E9}{\'e}             % é
\DeclareUnicodeCharacter{00F6}{\"o}             % ö
\DeclareUnicodeCharacter{00FC}{\"u}             % ü
\DeclareUnicodeCharacter{010C}{\v{C}}           % Č
\DeclareUnicodeCharacter{010D}{\v{c}}           % č
\DeclareUnicodeCharacter{0394}{\ensuremath{\Delta}}      % Δ
\DeclareUnicodeCharacter{03A3}{\ensuremath{\Sigma}}      % Σ
\DeclareUnicodeCharacter{03A9}{\ensuremath{\Omega}}      % Ω
\DeclareUnicodeCharacter{03B1}{\ensuremath{\alpha}}      % α
\DeclareUnicodeCharacter{03B2}{\ensuremath{\beta}}       % β
\DeclareUnicodeCharacter{03B3}{\ensuremath{\gamma}}      % γ
\DeclareUnicodeCharacter{03B4}{\ensuremath{\delta}}      % δ
\DeclareUnicodeCharacter{03B5}{\ensuremath{\varepsilon}} % ε
\DeclareUnicodeCharacter{03BA}{\ensuremath{\kappa}}      % κ
\DeclareUnicodeCharacter{03BB}{\ensuremath{\lambda}}     % λ
\DeclareUnicodeCharacter{03BC}{\ensuremath{\mu}}         % μ
\DeclareUnicodeCharacter{03BD}{\ensuremath{\nu}}         % ν
\DeclareUnicodeCharacter{03C0}{\ensuremath{\pi}}         % π
\DeclareUnicodeCharacter{03C1}{\ensuremath{\rho}}        % ρ
\DeclareUnicodeCharacter{03C3}{\ensuremath{\sigma}}      % σ
\DeclareUnicodeCharacter{03C4}{\ensuremath{\tau}}        % τ
\DeclareUnicodeCharacter{03C6}{\ensuremath{\varphi}}     % φ
\DeclareUnicodeCharacter{03C7}{\ensuremath{\chi}}        % χ
\DeclareUnicodeCharacter{03C8}{\ensuremath{\psi}}        % ψ
\DeclareUnicodeCharacter{03C9}{\ensuremath{\omega}}      % ω
\DeclareUnicodeCharacter{2013}{--}              % – (en dash)
\DeclareUnicodeCharacter{2014}{---}             % — (em dash)
\DeclareUnicodeCharacter{2018}{`}               % ‘
\DeclareUnicodeCharacter{2019}{'}               % ’
\DeclareUnicodeCharacter{201C}{``}              % “
\DeclareUnicodeCharacter{201D}{''}              % ”
\DeclareUnicodeCharacter{2070}{\ensuremath{^{0}}}        % ⁰
\DeclareUnicodeCharacter{2071}{\ensuremath{^{i}}}        % ⁱ
\DeclareUnicodeCharacter{2122}{\textsuperscript{TM}}     % ™
\DeclareUnicodeCharacter{2126}{\ensuremath{\Omega}}      % Ω (ohm sign)
\DeclareUnicodeCharacter{210F}{\ensuremath{\hbar}}       % ℏ
\DeclareUnicodeCharacter{2115}{\ensuremath{\mathbb{N}}}  % ℕ
\DeclareUnicodeCharacter{2102}{\ensuremath{\mathbb{C}}}  % ℂ
\DeclareUnicodeCharacter{211D}{\ensuremath{\mathbb{R}}}  % ℝ
\DeclareUnicodeCharacter{2124}{\ensuremath{\mathbb{Z}}}  % ℤ
\DeclareUnicodeCharacter{211A}{\ensuremath{\mathbb{Q}}}  % ℚ
\DeclareUnicodeCharacter{2192}{\ensuremath{\to}}         % →
\DeclareUnicodeCharacter{21D2}{\ensuremath{\Rightarrow}} % ⇒
\DeclareUnicodeCharacter{2200}{\ensuremath{\forall}}     % ∀
\DeclareUnicodeCharacter{2203}{\ensuremath{\exists}}     % ∃
\DeclareUnicodeCharacter{2208}{\ensuremath{\in}}         % ∈
\DeclareUnicodeCharacter{2209}{\ensuremath{\notin}}      % ∉
\DeclareUnicodeCharacter{2227}{\ensuremath{\wedge}}      % ∧
\DeclareUnicodeCharacter{2228}{\ensuremath{\vee}}        % ∨
\DeclareUnicodeCharacter{2229}{\ensuremath{\cap}}        % ∩
\DeclareUnicodeCharacter{222A}{\ensuremath{\cup}}        % ∪
\DeclareUnicodeCharacter{2245}{\ensuremath{\cong}}       % ≅
\DeclareUnicodeCharacter{2248}{\ensuremath{\approx}}     % ≈
\DeclareUnicodeCharacter{2260}{\ensuremath{\neq}}        % ≠
\DeclareUnicodeCharacter{2264}{\ensuremath{\leq}}        % ≤
\DeclareUnicodeCharacter{2265}{\ensuremath{\geq}}        % ≥
\DeclareUnicodeCharacter{22C5}{\ensuremath{\cdot}}       % ⋅
\DeclareUnicodeCharacter{2295}{\ensuremath{\oplus}}      % ⊕
\DeclareUnicodeCharacter{2297}{\ensuremath{\otimes}}     % ⊗

% ---- Mathematics & symbols ------------------------------------------
\usepackage{amsmath,amssymb,bm,mathrsfs,mathtools}
\usepackage{amsthm}             % theorem environments (load BEFORE hyperref)

% ---- Theorem environments (shared counter set, TECT canonical) -------
% All theorems / lemmas / propositions / corollaries / remarks share a
% single counter so cross-references read consistently across a paper.
% Definitions get a separate counter (definitions are numbered in their
% own sequence by editorial convention).
\newtheorem{theorem}{Theorem}
\newtheorem{lemma}[theorem]{Lemma}
\newtheorem{proposition}[theorem]{Proposition}
\newtheorem{corollary}[theorem]{Corollary}

\theoremstyle{remark}
\newtheorem{remark}[theorem]{Remark}
\newtheorem{example}[theorem]{Example}

\theoremstyle{definition}
\newtheorem{definition}[theorem]{Definition}
\newtheorem{conjecture}[theorem]{Conjecture}
\newtheorem{hypothesis}[theorem]{Hypothesis}

% ---- Tables / floats / graphics --------------------------------------
\usepackage{booktabs}           % \toprule / \midrule / \bottomrule
\usepackage{array}
\usepackage{graphicx}
\usepackage{enumitem}           % \begin{itemize}[leftmargin=...] options
\usepackage{hyphenat}           % \hyp{} for breakable hyphens in \texttt
\usepackage{seqsplit}           % \seqsplit{} for long unbreakable strings
\sloppy                         % allow stretchier inter-word spacing to
                                % reduce overfull \hbox warnings on long URLs
                                % and long \texttt tokens (paragraph-level).
\emergencystretch=3em           % last-resort hbox stretch budget

% ---- Cross-references (hyperref last among the standard packages,
%      cleveref AFTER hyperref) ----------------------------------------
\usepackage[%
  colorlinks=true,%
  linkcolor=blue,%
  citecolor=blue,%
  urlcolor=blue,%
  breaklinks=true,%
  bookmarks=true,%
  bookmarksnumbered=true%
]{hyperref}
\usepackage[capitalise,nameinlink]{cleveref}

% ---- TECT-canonical author block macros -----------------------------
% (defined here so every paper has the same author / affiliation style)
\newcommand{\tectauthor}{%
  \author{Jusang Lee}%
  \email{jtkor@outlook.com}%
  \affiliation{Independent researcher, Republic of Korea}%
}

% ---- TECT-canonical archive footer ----------------------------------
\newcommand{\tectarchivefooter}{%
  \par\medskip
  \noindent\small\textit{Canonical archive}:
  \url{https://github.com/TECT-OpenPhysics/TECT}.\par%
}

% ---- TECT TODO / AUDIT-FLAG / HONEST-SCOPE inline marks --------------
% Used in drafts to highlight pending audit items without disturbing
% the body text style. Suppress via \tectsuppressmarks (define empty).
\providecommand{\tectauditflag}[1]{\textcolor{red}{\textbf{[AUDIT-FLAG: #1]}}}
\providecommand{\tecttodo}[1]{\textcolor{orange}{\textbf{[TODO: #1]}}}
\providecommand{\tecthonestscope}[1]{\textit{Honest scope: #1.}}

% =====================================================================
% End of TECT canonical preamble
% =====================================================================

%   \begin{document}
%   ...
%   \end{document}
%
% Build (every paper): `pdflatex Paper-NN && pdflatex Paper-NN`
% (the second pass resolves \ref{} / \cref{} forward references).
% A bash/PowerShell wrapper that automates this is at
% Docs/papers/papers/_shared/build-paper.sh / build-paper.ps1.
% =====================================================================

% ---- Document class --------------------------------------------------
% PRD class option of revtex4-2 with [reprint,onecolumn] gives the same
% APS heading / by-line / abstract style as PRL but in a wider single
% column that handles long display equations and theorem environments
% without column-break artefacts. nofootinbib keeps citations out of
% footnotes. The 11pt option restores standard font metrics; with T1
% fontenc + lmodern this gives non-bitmap glyphs.
\documentclass[%
  aps,prd,reprint,onecolumn,nofootinbib,11pt,%
  amsmath,amssymb,floatfix,longbibliography%
]{revtex4-2}

% ---- Encoding & fonts ------------------------------------------------
\usepackage[T1]{fontenc}        % proper 8-bit font encoding (no font corruption)
\usepackage[utf8]{inputenc}     % allow UTF-8 source
\usepackage{lmodern}            % scalable Latin Modern (replaces bitmap CM)
\usepackage{textcomp}           % extra text symbols
\usepackage{microtype}          % subtle kerning improvements

% ---- UTF-8 → LaTeX character mappings --------------------------------
% Maps the Unicode characters that occur in TECT body text (legacy from
% the chat-content auto-archival rule, CLAUDE.md §4) to LaTeX-safe
% commands. Without these declarations, T1+inputenc rejects the chars
% with "Unicode character X not set up for use with LaTeX".
% Adding declarations centrally (here) is preferable to per-file edits.
\DeclareUnicodeCharacter{00A7}{\S}              % §  (section sign)
\DeclareUnicodeCharacter{00B2}{\ensuremath{^{2}}}        % ²
\DeclareUnicodeCharacter{00B3}{\ensuremath{^{3}}}        % ³
\DeclareUnicodeCharacter{00B5}{\ensuremath{\mu}}         % µ (micro)
\DeclareUnicodeCharacter{00B7}{\ensuremath{\cdot}}       % · (middle dot)
\DeclareUnicodeCharacter{00D7}{\ensuremath{\times}}      % ×
\DeclareUnicodeCharacter{00E4}{\"a}             % ä
\DeclareUnicodeCharacter{00E9}{\'e}             % é
\DeclareUnicodeCharacter{00F6}{\"o}             % ö
\DeclareUnicodeCharacter{00FC}{\"u}             % ü
\DeclareUnicodeCharacter{010C}{\v{C}}           % Č
\DeclareUnicodeCharacter{010D}{\v{c}}           % č
\DeclareUnicodeCharacter{0394}{\ensuremath{\Delta}}      % Δ
\DeclareUnicodeCharacter{03A3}{\ensuremath{\Sigma}}      % Σ
\DeclareUnicodeCharacter{03A9}{\ensuremath{\Omega}}      % Ω
\DeclareUnicodeCharacter{03B1}{\ensuremath{\alpha}}      % α
\DeclareUnicodeCharacter{03B2}{\ensuremath{\beta}}       % β
\DeclareUnicodeCharacter{03B3}{\ensuremath{\gamma}}      % γ
\DeclareUnicodeCharacter{03B4}{\ensuremath{\delta}}      % δ
\DeclareUnicodeCharacter{03B5}{\ensuremath{\varepsilon}} % ε
\DeclareUnicodeCharacter{03BA}{\ensuremath{\kappa}}      % κ
\DeclareUnicodeCharacter{03BB}{\ensuremath{\lambda}}     % λ
\DeclareUnicodeCharacter{03BC}{\ensuremath{\mu}}         % μ
\DeclareUnicodeCharacter{03BD}{\ensuremath{\nu}}         % ν
\DeclareUnicodeCharacter{03C0}{\ensuremath{\pi}}         % π
\DeclareUnicodeCharacter{03C1}{\ensuremath{\rho}}        % ρ
\DeclareUnicodeCharacter{03C3}{\ensuremath{\sigma}}      % σ
\DeclareUnicodeCharacter{03C4}{\ensuremath{\tau}}        % τ
\DeclareUnicodeCharacter{03C6}{\ensuremath{\varphi}}     % φ
\DeclareUnicodeCharacter{03C7}{\ensuremath{\chi}}        % χ
\DeclareUnicodeCharacter{03C8}{\ensuremath{\psi}}        % ψ
\DeclareUnicodeCharacter{03C9}{\ensuremath{\omega}}      % ω
\DeclareUnicodeCharacter{2013}{--}              % – (en dash)
\DeclareUnicodeCharacter{2014}{---}             % — (em dash)
\DeclareUnicodeCharacter{2018}{`}               % ‘
\DeclareUnicodeCharacter{2019}{'}               % ’
\DeclareUnicodeCharacter{201C}{``}              % “
\DeclareUnicodeCharacter{201D}{''}              % ”
\DeclareUnicodeCharacter{2070}{\ensuremath{^{0}}}        % ⁰
\DeclareUnicodeCharacter{2071}{\ensuremath{^{i}}}        % ⁱ
\DeclareUnicodeCharacter{2122}{\textsuperscript{TM}}     % ™
\DeclareUnicodeCharacter{2126}{\ensuremath{\Omega}}      % Ω (ohm sign)
\DeclareUnicodeCharacter{210F}{\ensuremath{\hbar}}       % ℏ
\DeclareUnicodeCharacter{2115}{\ensuremath{\mathbb{N}}}  % ℕ
\DeclareUnicodeCharacter{2102}{\ensuremath{\mathbb{C}}}  % ℂ
\DeclareUnicodeCharacter{211D}{\ensuremath{\mathbb{R}}}  % ℝ
\DeclareUnicodeCharacter{2124}{\ensuremath{\mathbb{Z}}}  % ℤ
\DeclareUnicodeCharacter{211A}{\ensuremath{\mathbb{Q}}}  % ℚ
\DeclareUnicodeCharacter{2192}{\ensuremath{\to}}         % →
\DeclareUnicodeCharacter{21D2}{\ensuremath{\Rightarrow}} % ⇒
\DeclareUnicodeCharacter{2200}{\ensuremath{\forall}}     % ∀
\DeclareUnicodeCharacter{2203}{\ensuremath{\exists}}     % ∃
\DeclareUnicodeCharacter{2208}{\ensuremath{\in}}         % ∈
\DeclareUnicodeCharacter{2209}{\ensuremath{\notin}}      % ∉
\DeclareUnicodeCharacter{2227}{\ensuremath{\wedge}}      % ∧
\DeclareUnicodeCharacter{2228}{\ensuremath{\vee}}        % ∨
\DeclareUnicodeCharacter{2229}{\ensuremath{\cap}}        % ∩
\DeclareUnicodeCharacter{222A}{\ensuremath{\cup}}        % ∪
\DeclareUnicodeCharacter{2245}{\ensuremath{\cong}}       % ≅
\DeclareUnicodeCharacter{2248}{\ensuremath{\approx}}     % ≈
\DeclareUnicodeCharacter{2260}{\ensuremath{\neq}}        % ≠
\DeclareUnicodeCharacter{2264}{\ensuremath{\leq}}        % ≤
\DeclareUnicodeCharacter{2265}{\ensuremath{\geq}}        % ≥
\DeclareUnicodeCharacter{22C5}{\ensuremath{\cdot}}       % ⋅
\DeclareUnicodeCharacter{2295}{\ensuremath{\oplus}}      % ⊕
\DeclareUnicodeCharacter{2297}{\ensuremath{\otimes}}     % ⊗

% ---- Mathematics & symbols ------------------------------------------
\usepackage{amsmath,amssymb,bm,mathrsfs,mathtools}
\usepackage{amsthm}             % theorem environments (load BEFORE hyperref)

% ---- Theorem environments (shared counter set, TECT canonical) -------
% All theorems / lemmas / propositions / corollaries / remarks share a
% single counter so cross-references read consistently across a paper.
% Definitions get a separate counter (definitions are numbered in their
% own sequence by editorial convention).
\newtheorem{theorem}{Theorem}
\newtheorem{lemma}[theorem]{Lemma}
\newtheorem{proposition}[theorem]{Proposition}
\newtheorem{corollary}[theorem]{Corollary}

\theoremstyle{remark}
\newtheorem{remark}[theorem]{Remark}
\newtheorem{example}[theorem]{Example}

\theoremstyle{definition}
\newtheorem{definition}[theorem]{Definition}
\newtheorem{conjecture}[theorem]{Conjecture}
\newtheorem{hypothesis}[theorem]{Hypothesis}

% ---- Tables / floats / graphics --------------------------------------
\usepackage{booktabs}           % \toprule / \midrule / \bottomrule
\usepackage{array}
\usepackage{graphicx}
\usepackage{enumitem}           % \begin{itemize}[leftmargin=...] options
\usepackage{hyphenat}           % \hyp{} for breakable hyphens in \texttt
\usepackage{seqsplit}           % \seqsplit{} for long unbreakable strings
\sloppy                         % allow stretchier inter-word spacing to
                                % reduce overfull \hbox warnings on long URLs
                                % and long \texttt tokens (paragraph-level).
\emergencystretch=3em           % last-resort hbox stretch budget

% ---- Cross-references (hyperref last among the standard packages,
%      cleveref AFTER hyperref) ----------------------------------------
\usepackage[%
  colorlinks=true,%
  linkcolor=blue,%
  citecolor=blue,%
  urlcolor=blue,%
  breaklinks=true,%
  bookmarks=true,%
  bookmarksnumbered=true%
]{hyperref}
\usepackage[capitalise,nameinlink]{cleveref}

% ---- TECT-canonical author block macros -----------------------------
% (defined here so every paper has the same author / affiliation style)
\newcommand{\tectauthor}{%
  \author{Jusang Lee}%
  \email{jtkor@outlook.com}%
  \affiliation{Independent researcher, Republic of Korea}%
}

% ---- TECT-canonical archive footer ----------------------------------
\newcommand{\tectarchivefooter}{%
  \par\medskip
  \noindent\small\textit{Canonical archive}:
  \url{https://github.com/TECT-OpenPhysics/TECT}.\par%
}

% ---- TECT TODO / AUDIT-FLAG / HONEST-SCOPE inline marks --------------
% Used in drafts to highlight pending audit items without disturbing
% the body text style. Suppress via \tectsuppressmarks (define empty).
\providecommand{\tectauditflag}[1]{\textcolor{red}{\textbf{[AUDIT-FLAG: #1]}}}
\providecommand{\tecttodo}[1]{\textcolor{orange}{\textbf{[TODO: #1]}}}
\providecommand{\tecthonestscope}[1]{\textit{Honest scope: #1.}}

% =====================================================================
% End of TECT canonical preamble
% =====================================================================


\begin{document}

\title{SO(10) anomaly cancellation via the spinor-trace method:
five Standard-Model perturbative anomaly coefficients $+$ one mixed
gravitational coefficient are symbolically exact zero per generation
(auxiliary anomaly-cancellation note)}

\author{Jusang Lee}
\email{jtkor@outlook.com}
\affiliation{Independent researcher, Republic of Korea}

\date{\today}

\begin{abstract}
Within the SO(10) grand-unification framework as employed by the
Topological Energy Condensate Theory (TECT) Pillar~4 candidate
embedding (Math157), we audit the per-generation cancellation of all
Standard-Model perturbative gauge anomalies and of the mixed
gravitational $\mathrm{U}(1)_Y$ anomaly. The five perturbative
coefficients
\[
  \mathrm{Tr}(\mathrm{U}(1)_Y^3),\quad
  \mathrm{Tr}(\mathrm{SU}(2)_L^3),\quad
  \mathrm{Tr}(\mathrm{SU}(3)_c^3),\quad
  \mathrm{Tr}(\mathrm{U}(1)_Y\cdot \mathrm{SU}(2)_L^2),\quad
  \mathrm{Tr}(\mathrm{U}(1)_Y\cdot \mathrm{SU}(3)_c^2)
\]
and the one gravitational mixed coefficient
$\mathrm{Tr}(\mathrm{U}(1)_Y)$ vanish \emph{symbolically exactly} per
generation when summed over the chiral content of the SO(10)
$\mathbf{16}$ spinor reduced under
$\mathrm{SU}(5)\times\mathrm{U}(1)_\chi
\to\mathrm{SU}(3)_c\times\mathrm{SU}(2)_L\times\mathrm{U}(1)_Y$. The
SU(5) sub-content $\overline{\mathbf{5}}\oplus\mathbf{10}$ already
saturates all six perturbative + gravitational SM anomalies; the
SO(10) singlet $\mathbf{1}(5)$ (the right-handed neutrino, hereafter
RHN) carries hypercharge $Y=0$ and therefore contributes \emph{zero}
to any $\mathrm{U}(1)_Y$ trace. The RHN's genuine consistency role is
(i) closure of the Witten $\mathrm{SU}(2)$ global anomaly
(non-perturbative mod-$2$, \emph{not} part of the perturbative
six-coefficient list); and (ii) when $\mathrm{U}(1)_{B-L}$ is gauged
in the SO(10) extension, saturation of the mixed gravitational
$\mathrm{U}(1)_{B-L}$ anomaly. The result is conditional on (a) the
SO(10) embedding of TECT Pillar~4 (Math157) and (b) the standard
SU(5) $\to$ SM hypercharge assignment via
$Y = T_{3R}+\tfrac{1}{2}(B-L)$. This paper is an independent
auxiliary anomaly-cancellation note and is \emph{not} coupled to the
unrelated three-generation question (Paper~07). Math319 V1
verification target adopts these six cancellations as its
symbolic-anomaly check.
\end{abstract}

\maketitle

% =====================================================================
\section{Introduction}\label{sec:intro}
% =====================================================================
Anomaly cancellation is a cornerstone of consistent quantum field
theory. For a four-dimensional chiral gauge theory with gauge group
$G_{\mathrm{SM}}=\mathrm{SU}(3)_c\times\mathrm{SU}(2)_L\times
\mathrm{U}(1)_Y$, perturbative consistency requires the vanishing of
five chiral triangle anomaly coefficients,
\begin{equation}\label{eq:five-pert}
  \mathrm{Tr}(\mathrm{U}(1)_Y^3),\;
  \mathrm{Tr}(\mathrm{SU}(2)_L^3),\;
  \mathrm{Tr}(\mathrm{SU}(3)_c^3),\;
  \mathrm{Tr}(\mathrm{U}(1)_Y\cdot \mathrm{SU}(2)_L^2),\;
  \mathrm{Tr}(\mathrm{U}(1)_Y\cdot \mathrm{SU}(3)_c^2),
\end{equation}
together with the vanishing of the mixed gravitational anomaly
\begin{equation}\label{eq:grav}
  \mathrm{Tr}(\mathrm{U}(1)_Y),
\end{equation}
which is the coefficient of the graviton--graviton--$\mathrm{U}(1)_Y$
triangle. The traces in \eqref{eq:five-pert}--\eqref{eq:grav} are
sums over all left-handed (or, equivalently, all $L$-and-$R$-
conjugate) chiral fermions of the theory.

Two consistency conditions of the SM are not entries in
\eqref{eq:five-pert}--\eqref{eq:grav}:
\begin{itemize}
  \item the \emph{Witten $\mathrm{SU}(2)_L$ global anomaly}
    \cite{Witten1982}, which is a non-perturbative
    $\mathbb{Z}_2$ statement requiring an even number of chiral
    $\mathrm{SU}(2)_L$ doublets;
  \item the gauged $\mathrm{U}(1)_{B-L}$ mixed gravitational anomaly
    $\mathrm{Tr}(B-L)$, which appears \emph{only} when $B-L$ is
    promoted to a gauge symmetry (as in the SO(10) extension of the
    SM).
\end{itemize}
Both are independent of \eqref{eq:five-pert}--\eqref{eq:grav}; their
treatment in TECT is recorded below as a separate audit
(\S\ref{sec:rhn-genuine-role}).

The SO(10) grand-unification group naturally houses the SM gauge
structure. The spinor representation $\mathbf{16}$ of SO(10) contains
one complete SM family plus one electroweak singlet (the RHN).
Math157 exploits this to derive all SM gauge anomaly coefficients from
the group-theoretic trace formulas of SO(10).

% =====================================================================
\section{Setting and main result}\label{sec:setup}
% =====================================================================
The SO(10) spinor $\mathbf{16}$ decomposes under
$\mathrm{SU}(5)\times\mathrm{U}(1)_\chi$ as
\begin{equation}\label{eq:16dec}
  \mathbf{16} \;=\; \mathbf{10}(1) \,\oplus\, \overline{\mathbf{5}}(-3) \,\oplus\, \mathbf{1}(5),
\end{equation}
and further to the SM gauge group as
\begin{align}
  \mathbf{10}(1) &\;\to\;
    \underbrace{(\mathbf{3},\mathbf{2})_{1/6}}_{Q_L}
    \;\oplus\;
    \underbrace{(\overline{\mathbf{3}},\mathbf{1})_{-2/3}}_{u_R^c}
    \;\oplus\;
    \underbrace{(\mathbf{1},\mathbf{1})_{1}}_{e_R^c}, \label{eq:10dec}\\
  \overline{\mathbf{5}}(-3) &\;\to\;
    \underbrace{(\mathbf{1},\mathbf{2})_{-1/2}}_{L_L}
    \;\oplus\;
    \underbrace{(\overline{\mathbf{3}},\mathbf{1})_{1/3}}_{d_R^c}, \label{eq:5bardec}\\
  \mathbf{1}(5) &\;\to\;
    \underbrace{(\mathbf{1},\mathbf{1})_{0}}_{\nu_R^c} \;\;\;
    \text{(the RHN; SM hypercharge $Y=0$).} \label{eq:singletdec}
\end{align}
The hypercharge assignment in \eqref{eq:10dec}--\eqref{eq:singletdec}
is the standard $\mathrm{SU}(5)$-derived
$Y = T_{3R} + \tfrac{1}{2}(B-L)$.

\begin{theorem}[SO(10) anomaly cancellation via the spinor-trace
method]\label{thm:main}
Under hypotheses
\begin{itemize}
  \item[(\textbf{H1})] the SO(10) embedding (Math157) expressing the
    SM gauge group and one generation of SM matter as the chiral
    content of the SO(10) $\mathbf{16}$ spinor as in
    \eqref{eq:16dec}--\eqref{eq:singletdec};
  \item[(\textbf{H2})] the standard $\mathrm{SU}(5)$-derived
    hypercharge assignment $Y=T_{3R}+\tfrac{1}{2}(B-L)$ on each
    irreducible component (Math157-AddD),
\end{itemize}
the five perturbative SM gauge anomaly coefficients of
\eqref{eq:five-pert} and the one mixed gravitational coefficient
\eqref{eq:grav} are \emph{symbolically exact zero} per generation:
\begin{equation}\label{eq:six-zero}
\boxed{\;
\begin{array}{ll}
  \mathrm{Tr}(\mathrm{U}(1)_Y^3) = 0, &
  \mathrm{Tr}(\mathrm{SU}(2)_L^3) = 0, \\
  \mathrm{Tr}(\mathrm{SU}(3)_c^3) = 0, &
  \mathrm{Tr}(\mathrm{U}(1)_Y \cdot \mathrm{SU}(2)_L^2) = 0, \\
  \mathrm{Tr}(\mathrm{U}(1)_Y \cdot \mathrm{SU}(3)_c^2) = 0, &
  \mathrm{Tr}(\mathrm{U}(1)_Y) = 0,
\end{array}\;}
\end{equation}
where the trace is over the chiral content of \emph{one} $\mathbf{16}$
of SO(10) (i.e.\ one SM generation $+$ one RHN). The RHN $\mathbf{1}(5)$
contributes \emph{zero} to every line of \eqref{eq:six-zero} because
its SM hypercharge is $Y=0$.

The result is \textbf{T6 \textsc{proved conditional}} on (H1), (H2),
and the textbook-standard SU(5) representation theory employed in
\eqref{eq:10dec}--\eqref{eq:5bardec}.
\end{theorem}

\begin{remark}[The five-coefficient sub-saturation by
$\overline{\mathbf{5}}\oplus\mathbf{10}$]\label{rmk:5bar10sat}
The five perturbative coefficients of \eqref{eq:five-pert} are
\emph{already} saturated by the SU(5) sub-content
$\overline{\mathbf{5}}\oplus\mathbf{10}$ alone: this is a
classical result of Georgi--Glashow SU(5) GUT phenomenology
\cite{GeorgiGlashow1974}. The RHN $\mathbf{1}(5)$ is unnecessary for
\eqref{eq:five-pert}. Its role in
\eqref{eq:six-zero} is via $\mathrm{Tr}(\mathrm{U}(1)_Y)$; but since
the RHN has $Y=0$, it contributes zero there as well, and the
gravitational mixed anomaly is also saturated by
$\overline{\mathbf{5}}\oplus\mathbf{10}$ alone.
\end{remark}

\begin{remark}[Multi-generation extension]\label{rmk:multigen}
For $N_g$ generations the same theorem holds with all coefficients
in \eqref{eq:six-zero} multiplied by $N_g$. Cancellation is therefore
$N_g$-independent. Anomaly cancellation \emph{does not} fix the value
of $N_g$. (See \S\ref{sec:rhn-genuine-role} on the separate Witten
SU(2) global anomaly, which constrains $N_g\bmod 2$ only via the
total chiral $\mathrm{SU}(2)_L$ doublet count.)
\end{remark}

% =====================================================================
\section{Proof: term-by-term spinor trace}\label{sec:proof}
% =====================================================================
\paragraph{Step 1 (SO(10) spinor decomposition).}
The SO(10) Clifford algebra and the $\mathbf{16}$ spinor
representation are computed via the standard $32\times 32$ Dirac
gamma-matrix construction of \cite{GeorgiBook}; the
$\mathrm{SU}(5)\times \mathrm{U}(1)_\chi$ branching
\eqref{eq:16dec} is read off from the SO(10) Dynkin diagram via the
extended-Dynkin folding \cite{Slansky1981}. The further reduction to
$\mathrm{SU}(3)_c\times\mathrm{SU}(2)_L\times\mathrm{U}(1)_Y$ in
\eqref{eq:10dec}--\eqref{eq:5bardec} follows the standard
SU(5) $\to$ SM branching (Math157, \cite{Math157}).

\paragraph{Step 2 (SM hypercharge assignment).}
Each chiral component receives its SM hypercharge via
$Y = T_{3R} + \tfrac{1}{2}(B-L)$, where $T_{3R}$ is the third
component of right-handed isospin (read from the SO(10)
$\mathrm{SU}(2)_R$ sub-algebra) and $(B-L)$ is the baryon-minus-lepton
charge (Math157-AddD, \cite{Math157AddD}).

\paragraph{Step 3 (trace formulas).}
For each anomaly coefficient, the trace is computed by summing the
appropriate Dynkin index and charge over all left-handed chiral
fermions in the spinor:
\begin{equation}\label{eq:trace-formula}
  \mathrm{Tr}(G^a G^b G^c) \;=\; \sum_{f\in\mathbf{16}} d^{abc}_{R_f}\,n_f,
\end{equation}
where $d^{abc}_{R_f}$ is the symmetric anomaly coefficient of the
representation $R_f$ of the chiral fermion $f$ under generators
$G^a,G^b,G^c$, and $n_f=1$ for left-handed fermions, $n_f=-1$ for
right-handed (here we use the $L$-conjugate convention so all $n_f=1$).

\paragraph{Step 4 (term-by-term cancellation).}
The five perturbative coefficients of \eqref{eq:five-pert} evaluate
on the chiral content of $\overline{\mathbf{5}}\oplus\mathbf{10}$ to:
\begin{align}
  \mathrm{Tr}(\mathrm{U}(1)_Y^3) &=
    3 \cdot 2 \cdot (\tfrac{1}{6})^3
    + 3\cdot(-\tfrac{2}{3})^3
    + (1)^3
    + 2\cdot(-\tfrac{1}{2})^3
    + 3\cdot(\tfrac{1}{3})^3 = 0, \\
  \mathrm{Tr}(\mathrm{SU}(2)_L^3) &= 0
    \quad\text{(symmetric anomaly trivially vanishes for SU(2);}\\
  &\quad\text{Witten global anomaly is a separate mod-$2$ statement)}, \nonumber\\
  \mathrm{Tr}(\mathrm{SU}(3)_c^3) &=
    2\cdot 1 + (-1) + (-1) = 0
    \quad\text{(net triangle vanishes)}, \\
  \mathrm{Tr}(\mathrm{U}(1)_Y \cdot \mathrm{SU}(2)_L^2) &=
    3\cdot \tfrac{1}{6} \cdot \tfrac{1}{2}
    + (-\tfrac{1}{2}) \cdot \tfrac{1}{2} = 0, \\
  \mathrm{Tr}(\mathrm{U}(1)_Y \cdot \mathrm{SU}(3)_c^2) &=
    2 \cdot \tfrac{1}{6} \cdot \tfrac{1}{2}
    + (-\tfrac{2}{3})\cdot \tfrac{1}{2}
    + \tfrac{1}{3}\cdot \tfrac{1}{2} = 0,
\end{align}
where the SU(2) (resp.\ SU(3)) Dynkin index $\tfrac{1}{2}$ for the
fundamental, and multiplicities follow from
\eqref{eq:10dec}--\eqref{eq:5bardec}. The mixed gravitational
coefficient evaluates to
\begin{equation}\label{eq:trY-zero}
  \mathrm{Tr}(\mathrm{U}(1)_Y) =
    3\cdot 2 \cdot \tfrac{1}{6}
    + 3\cdot(-\tfrac{2}{3})
    + 1
    + 2\cdot(-\tfrac{1}{2})
    + 3\cdot \tfrac{1}{3} = 0,
\end{equation}
again from $\overline{\mathbf{5}}\oplus\mathbf{10}$ alone. The RHN
contributes zero everywhere (its SU(2)$_L$, SU(3)$_c$, and $Y$
charges are all zero). \qed

% =====================================================================
\section{The RHN's genuine consistency role}\label{sec:rhn-genuine-role}
% =====================================================================
The earlier draft of this paper (Rev 3) asserted that the RHN cancels
the $\mathrm{U}(1)_Y$ anomaly of the $\overline{\mathbf{5}}(-3)$
representation. As recorded in \S\ref{sec:proof} and
Remark~\ref{rmk:5bar10sat}, this statement is incorrect: the RHN has
SM hypercharge $Y=0$ and contributes zero to every $\mathrm{U}(1)_Y$
trace, while $\overline{\mathbf{5}}\oplus\mathbf{10}$ already saturates
all six SM perturbative + gravitational anomalies of
\eqref{eq:six-zero}.

The RHN does play a real role in three independent consistency audits
of the SO(10) extension:
\begin{enumerate}
  \item \textbf{Witten $\mathrm{SU}(2)_L$ global anomaly}
    \cite{Witten1982}. This is a non-perturbative
    $\mathbb{Z}_2$ statement: a four-dimensional theory with an odd
    number of chiral $\mathrm{SU}(2)_L$ doublets is inconsistent at
    the path-integral level. The SM has, per generation, the doublets
    $Q_L$ (with multiplicity~$3$ from colour) and $L_L$
    (multiplicity $1$), giving a total of $4$ chiral doublets per
    generation, which is even; the global anomaly already cancels
    \emph{without} the RHN. The RHN, being an SU(2) singlet,
    contributes zero doublets and therefore plays no role in the
    Witten cancellation either. Its consistency role here is
    \emph{neutral} (it does not break the cancellation; equivalently,
    SU(2) anomaly cancellation does not require it).

  \item \textbf{Mixed gravitational $\mathrm{U}(1)_{B-L}$ anomaly}.
    When $B-L$ is gauged in the SO(10) extension (the
    $\mathrm{U}(1)_{B-L}$ subgroup of the SO(10)
    Pati--Salam $\mathrm{SU}(4)_C$), the trace
    $\mathrm{Tr}(B-L)$ is a new mixed gravitational anomaly. The RHN
    has $(B-L) = +1$ (lepton number of opposite sign because it is the
    $L$-conjugate of $\nu_R$); it is required to make
    $\mathrm{Tr}(B-L) = 0$ per generation. This is the RHN's strongest
    classical role.

  \item \textbf{SO(10) algebra anomaly freedom}. The Lie algebra
    $\mathfrak{so}(10)$ is anomaly-free at the simple-algebra level
    (\cite{Slansky1981}, Sec.\,4); its representations admit the
    spinor $\mathbf{16}$ as the smallest faithful spinor that contains
    one full SM generation. The RHN appears as the singlet component
    of $\mathbf{16}$; its presence is the consequence, not the cause,
    of completing the spinor.
\end{enumerate}
None of (1)--(3) is identified with the perturbative-gauge cancellation
\eqref{eq:five-pert} or with the gravitational $\mathrm{Tr}(Y)$
cancellation \eqref{eq:trY-zero}. The corrected statement is therefore:
\emph{the RHN is required for the $B-L$-gauged extension and for
completion of the $\mathbf{16}$ spinor; it is not required for SM
perturbative or gravitational $\mathrm{U}(1)_Y$ anomaly cancellation.}

% =====================================================================
\section{Numerical verification}\label{sec:num}
% =====================================================================
The six trace identities \eqref{eq:six-zero} were cross-checked by
symbolic algebra over the SU(5) representation content
\eqref{eq:10dec}--\eqref{eq:5bardec}; each line of \eqref{eq:six-zero}
returns the identically-zero rational number, not a finite-precision
estimate. A subsequent numerical evaluation in machine arithmetic
returns no residual at machine precision; this is a redundant
cross-check and \emph{not} a numerical estimate.

\begin{remark}[Symbolic versus numerical exactness]\label{rmk:symbolic}
The earlier draft phrasing ``numerically to six-digit precision'' was
imprecise. The cancellation \eqref{eq:six-zero} is a symbolic identity
in the rational numbers; the appropriate phrasing is ``symbolically
exact zero, numerically cross-checked''. The Math319 V1 verification
target adopts the symbolic form.
\end{remark}

% =====================================================================
\section{Discussion and outlook}\label{sec:discussion}
% =====================================================================
The anomaly-cancellation result \eqref{eq:six-zero} is a non-trivial
consistency check of the SO(10) embedding derived from the BCC
condensate (Math157, Math229). It is independent of the
three-generation question (Paper~07): cancellation holds per
generation, for any number $N_g$ of generations. Anomaly cancellation
therefore does \emph{not} fix the SM generation count; that question
remains the canonical Pillar~7 \textbf{T1 OPEN} problem.

The RHN's classical role is correctly identified, post-correction, as
$B-L$ anomaly closure and SO(10) spinor completion, not as
$\mathrm{U}(1)_Y$ saturation (which is entirely absorbed by the
$\overline{\mathbf{5}}\oplus\mathbf{10}$ sub-content of the SU(5)
GUT). This paper is the auxiliary anomaly-cancellation note for the
TECT Pillar~4 SO(10) embedding programme; it is decoupled from
Paper~07 (three-generation forcing), which carries an independent
audit-flag and rewrite obligation.

% =====================================================================
\begin{thebibliography}{99}
\bibitem{Witten1982} E.~Witten,
  ``An $\mathrm{SU}(2)$ anomaly,''
  Phys.\ Lett.\ B \textbf{117}, 324 (1982).
\bibitem{GeorgiGlashow1974} H.~Georgi and S.~L.~Glashow,
  ``Unity of all elementary-particle forces,''
  Phys.\ Rev.\ Lett.\ \textbf{32}, 438 (1974).
\bibitem{GeorgiBook} H.~Georgi,
  \emph{Lie Algebras in Particle Physics}, 2nd ed.\ (Westview, 1999).
\bibitem{Slansky1981} R.~Slansky,
  ``Group theory for unified model building,''
  Phys.\ Rep.\ \textbf{79}, 1 (1981).
\bibitem{Math157} TECT internal note Math157,
  ``SO(10) spinor decomposition and TECT Pillar~4 embedding,''
  TECT canonical archive.
\bibitem{Math157AddD} TECT internal note Math157-AddD,
  ``RHN identification and $B-L$ assignment in the SO(10) embedding,''
  TECT canonical archive.
\end{thebibliography}

\begin{acknowledgments}
The author thanks Math157 and Math157-AddD for the trace-method
formalism and RHN identification, and the hostile-referee audit
(Math314-AddC, 2026-05-02) for catching the RHN--$\mathrm{U}(1)_Y$
mis-attribution corrected in Rev~4. Numerical cross-check was
performed with SymPy symbolic algebra. This work is part of TECT,
canonical archive at
\url{https://github.com/TECT-OpenPhysics/TECT}.
\end{acknowledgments}

\end{document}
